% L15a strategy deck: the Stone-Geary weighting rule.
% One short deck per rule on the strategy menu, read next to the rule's code file
% (strategies/weighting/stone_geary.jl). Written by make_strategy_decks.py (session scratchpad, 2026-11-27) and then
% kept by hand: edit this file, not the generator. Build with make in this folder: LaTeX runs from strategies/,
% so paths are relative to it.
\documentclass[aspectratio=169,11pt]{beamer}
\usepackage{vnslides}

\newcommand{\repo}{https://github.com/varnerlab/CHEME-5660-CourseRepository-Fall-2026/blob/main/lectures/week-15/L15a}

\title{\texorpdfstring{{\vnheavy L15a}}{L15a}: Stone-Geary}
\subtitle{Weighting Rule, CHEME 5660 Strategy Menu, Cornell University}
\author{Copyright Jeffrey Varner 2026. All Rights Reserved}

\begin{document}

\vntitlepage

\begin{frame}{Disclaimer and Risks}
  \vspace{0.6em}
  \textbf{\ink{Training only.}} This content is for training and informational
  purposes only. The teaching team does not make, give, or endorse any offer or
  solicitation to buy or sell securities or derivative products, or provide
  investment or trading advice or strategy.

  \vspace{1.1em}
  \textbf{\ink{Risk.}} Trading involves risk and is generally not recommended for
  those with limited resources, experience, or a low risk tolerance. Past
  performance does not guarantee future success.

  \vspace{1.1em}
  \textbf{\ink{Responsibility.}} You are responsible for your investment decisions
  based on your financial circumstances, objectives, risk tolerance, and
  liquidity needs.
\end{frame}

\begin{frame}{The Idea}
  \vspace{0.3em}
  The Stone-Geary utility gives every ticker in the basket a baseline of 1\% of the wealth, and splits the
  rest as in Cobb-Douglas:
  \begin{itemize}\setlength{\itemsep}{0.5em}
    \item \hilite{Long only}: the rule never sells a ticker short.
    \item \hilite{Baseline}: every ticker holds at least 1\% of the wealth, even one with a negative preference weight. The baseline keeps every picked ticker in the portfolio.
    \item \hilite{Signal}: above the baseline, a ticker with a higher preference weight gets more of the wealth.
  \end{itemize}

  \vspace{0.5em}
  The basket comes from the picking rule. Each preference weight is the day's value at the decision.
\end{frame}

\begin{frame}{The Preference Weight}
  \vspace{0.3em}
  At each close $t$, the engine scores every ticker $i$ with its preference weight $\gamma_{i}(t)$:
  {\setlength{\jot}{3pt}
  \begin{align*}
    \tilde{g}_{i,t} &= \alpha_{i}(t)+\beta_{i}(t)\,\tilde{g}_{M,t} && \text{the expected growth rate (1/yr)}\\
    \gamma_{i}(t) &= \tanh\left(h\,\frac{\tilde{g}_{i,t}}{\max\left(|\beta_{i}(t)|,\,\beta_{\text{floor}}\right)^{\xi_{t}}}\right)
      && h=1\text{ year},\ \beta_{\text{floor}}=0.1
  \end{align*}}%
  where the single index model gives $\alpha_{i}(t)$ and $\beta_{i}(t)$, $\tilde{g}_{M,t}$ is the 252-day
  exponential moving average of market growth, and the market state $\xi_{t}$ is positive when the market falls.

  \vspace{0.5em}
  The preference weight lies between $-1$ and $1$, with the sign of the expected growth rate.
\end{frame}

\begin{frame}{How the Market State Tilts the Preference Weight}
  \vspace{0.3em}
  The market state $\xi_{t}$ sets how the market sensitivity scales the expected growth rate:
  \begin{center}
    \includegraphics{slides/figs/Fig-L15a-Deck-Tilt.pdf}
  \end{center}
  At the same positive expected growth rate, the preference weight favors a low market sensitivity when the
  market falls and a high one when the market rises.
\end{frame}

\begin{frame}{The Stone-Geary Weights}
  \vspace{0.3em}
  Each of the $K$ tickers holds a baseline of $b_{i}$ shares, with $b_{i}S_{i}=0.01\,W$. The rest of the wealth,
  $(1-0.01K)\,W$, goes to the shares above the baseline, $n_{i}-b_{i}$:
  {\setlength{\jot}{2pt}
  \begin{align*}
    &\max_{n}\ \sum_{\gamma_{i}>0}\gamma_{i}\ln\left(n_{i}-b_{i}\right)\ \ \text{subject to}\ \ \sum_{i}n_{i}S_{i}=W
      && \text{utility and budget}\\
    &\left(n_{i}-b_{i}\right)S_{i}=\frac{\gamma_{i}}{\sum_{\gamma_{j}>0}\gamma_{j}}\left(1-0.01K\right)W
      && \text{Cobb-Douglas on the rest}\\
    \Longrightarrow\quad&\boxed{w_{i}=0.01+\left(1-0.01K\right)\frac{\max(\gamma_{i},0)}{\sum_{j}\max(\gamma_{j},0)}}
  \end{align*}}%
  For three tickers with $\gamma=(0.6,\ 0.3,\ -0.1)$, the weights are $(0.657,\ 0.333,\ 0.010)$. The baselines fit within the wealth only when $K\leq 100$.
\end{frame}

\begin{frame}{The Weights of Three Made-Up Tickers}
  \vspace{0.3em}
  The three tickers with $\gamma=(0.6,\ 0.3,\ -0.1)$ get these weights:
  \begin{center}
    \includegraphics{slides/figs/Fig-L15a-Deck-Weights-Stone-Geary.pdf}
  \end{center}
  Every ticker of the basket keeps its 1\% baseline, even the third ticker with a negative preference weight. On
  top of their baselines, the first two tickers split the other 97\% in the Cobb-Douglas proportions.
\end{frame}

\begin{frame}{Under the Trading Limits}
  \vspace{0.3em}
  At each decision, the engine chooses the share counts $n_{i}$ that maximize the Stone-Geary utility
  within the trading limits:
  \[
    U(n;\gamma)=\sum_{\gamma_{i}>0}\gamma_{i}\ln\left(n_{i}-b_{i}\right),\qquad n_{i}\geq b_{i}
  \]
  \begin{itemize}\setlength{\itemsep}{0.5em}
    \item \hilite{Shared limits}: the weight, volatility, and trading-cost limits of every strategy apply. With the hedge on, so does the market sensitivity limit, and the engine can buy up to five tickers with zero or negative market sensitivity from its hedge pool.
    \item \hilite{Baseline}: a ticker with $\gamma_{i}\leq 0$ holds its baseline $b_{i}$ shares.
  \end{itemize}

  \vspace{0.3em}
  The cash stays at or above zero. When a limit stops the engine from reaching the Stone-Geary weights, the engine holds the best portfolio within the limits, and cash can remain.
\end{frame}

\begin{frame}{Reading the Code}
  \vspace{0.3em}
  The rule's code file is
  \href{\repo/strategies/weighting/stone_geary.jl}{\texttt{strategies/weighting/stone\_geary.jl}}:
  \begin{itemize}\setlength{\itemsep}{0.5em}
    \item \texttt{weigh\_stone\_geary(gamma; baseline)} computes the weights before the trading limits. The engine does not call it.
    \item \texttt{MyL15aStoneGearyRule} is the type the engine calls. It holds the setting \texttt{baseline~=~0.01}.
    \item \texttt{MPC.allocate(rule, p)} passes the decision's problem \texttt{p} to the engine's \texttt{MPC.MyStoneGearyRule}, which maximizes the utility within the limits.
  \end{itemize}

  \vspace{0.5em}
  The \href{\repo/strategies/CHEME-5660-L15a-Strategy-Menu-Fall-2026.ipynb}{strategy menu} lists every rule, and the \href{\repo/engine/CHEME-5660-L15a-Trading-Engine-Fall-2026.ipynb}{trading engine} describes the shared limits.
\end{frame}

\end{document}
